\section{Outstanding analytical tasks}\label{sec:outstanding}

The transition-scale consumer is selected, and its outer reduction and
adaptive-reference comparison costs are paid. The remaining obstacle is
a value estimate at exceptional centers. The density and averaged-value
conclusions, including the quantitative short-window bounds, leave
that obstacle open. The following tasks distinguish
the missing arithmetic information from further verification and changes
of scope.

\paragraph{1. Acquire new arithmetic information at exceptional centers.}
For the unchanged band \(\lambda=c/L_X\), \(1/2\le c\le1\), and
\(P(X)=\lfloor L_X^2\rfloor\), estimate the complete
\(R_{P(X)}(X,c)\) in \eqref{eq:exact-relative-residual} strongly
enough for \eqref{eq:adaptive-allowance} to lie below \(m_*\).
Alternatively estimate the actual signed nonlinear logarithmic
contribution, retaining its amplitude and phase dependence. A new input
must specify a proof route and a complete error budget; repeating the
residual's definition or the desired signed inequality is not that input.
The field error \(X^{-1/4+o(1)}\) is already paid and does not bound
the arithmetic residual. The short-window Green bound allows total
mass \(O_\rho(\ell L_U^{-1/3})\), which grows on the admissible
windows. A persistence or maximum-to-average estimate would have to
control the actual trajectory \(X\mapsto(1/L_X,X)\), including its
time derivative contribution where differentiable.

\paragraph{2. Preserve the proved limits of the tested methods.}
The exact quadratic-Jensen recipe fails at arbitrarily high centers;
even its diagnostic ideal cosine subtraction fails by at least
\(1/1000\). Every fixed Euler reference has unbounded residuals,
and its specified signed tangent fails uniformly over regulators.
These statements do not decide a growing-reference tangent, spatially
varying tangent heights or all nonlinear signed methods. Dividing by
the reflected reference would meet its real zeros and requires a
separate small-value argument. No such quotient is assumed zero-free.
Keep these tested methods parked unless the proposed change supplies
information outside their proved failure class.

\paragraph{3. Resolve uniformity only in the domain actually needed.}
The density, finite-power and subcritical exponential mean conclusions
apply at \(c=1\), with \(\lambda=1/L_X\) varying between centers.
Theorem~\ref{thm:b10-averaged} uses a cutoff fixed over each translated
interval. Corollary~\ref{cor:b10-green} gives a rate for the Green mean
and the fixed-threshold exceptional proportion with that cutoff.
The varying-cutoff residual, the extension throughout \(c\in[1/2,1]\),
and the shrinking-depth regime need separate uniform estimates.
The fourth exponential mean is bounded and its convergence remains
open. The asymptotic statements begin at a finite unevaluated onset.

\paragraph{4. Retain the separate time-zero arithmetic problem.}
At time zero the parent target remains an upper bound for
\(\int K\log|F|-P+G(X,1)\) below the chosen defect margin.
The increasing-order Riemann--Siegel remainder in
\eqref{eq:RS-log-transfer} controls its logarithmic error but not the
retained arithmetic value. The shifted-zeta identity
\eqref{eq:shifted-zeta-divisor} retains an unknown zero potential
and signed difference; it does not settle this parent question.
The localized-correlation and collision-constraint tracks likewise
retain their separate data requirements and ownership.

\paragraph{5. Check the foundation and make constants effective when useful.}
The imported complete-field estimate, its fixed buffers and Gaussian
exterior, the source normalization, opposite heat clocks, divisor
signs and logarithmic tails retain their original internal-review
status. The B10 averaged theorem has a separate internal audit with
contributing agents and shared inputs. Specialist review remains
outstanding. The finite provider constants and onsets are unevaluated.

\paragraph{6. Use numerics for a specified analytical question.}
The new results are analytical. Existing numerical certificates retain
their finite supporting scope in Section~\ref{sec:evidence}. A future
computation should calibrate a named constant, test a specified proposed
estimate or inspect a proved failure mechanism, with a fixed comparator
and budget. A height scan cannot establish the required every-center
bound. Further execution should follow a new arithmetic estimate or a
concrete proof question; no additional task is issued merely to rename
the remaining residual.
